%── environments/story-blocks.tex ───────────────────────────────────────────
%
% Entornos auxiliares del poster.
% IMPORTANTE: todos usan StoryAccent y StoryAccentFaint (nombres resueltos
% con \colorlet en main-story.tex), nunca \StoryAccentColor directamente.
%───────────────────────────────────────────────────────────────────────────

% ─── Highlight inline de expresión matemática ────────────────────────────────
% Fondo muy sutil en el color de acento de la historia.
% Uso: la función \inlinemath{f(x;\theta)} parametriza la densidad.
\newcommand{\inlinemath}[1]{%
  \tcbox[
    enhanced,
    frame hidden,
    colback=StoryAccentFaint,
    arc=1.5pt,
    boxrule=0pt,
    left=1.5pt, right=1.5pt, top=0pt, bottom=0pt,
    nobeforeafter,
    box align=base,
  ]{\small$\displaystyle#1$}%
}

% ─── Dato clave ──────────────────────────────────────────────────────────────
% Bloque compacto con borde izquierdo, sin fondo.
% Uso: \KeyFact{La varianza estacionaria es $\sigma^2/(2\theta)$}
\newcommand{\KeyFact}[1]{%
  \begin{tcolorbox}[
    leftaccent=StoryAccent,
    fontupper=\fontsize{7}{9}\selectfont\bfseries\color{InkBlack},
  ]
    #1
  \end{tcolorbox}%
}

% ─── Separador con label opcional ────────────────────────────────────────────
% Uso: \StorySep  o  \StorySep{Derivación}
\newcommand{\StorySep}[1][]{%
  \par\vspace{3pt}%
  \GrayRule%
  \ifx&#1&%
  \else
    \par\vspace{1pt}%
    {\fontsize{5.5}{6}\selectfont\bfseries\color{InkLight}\MakeUppercase{#1}}%
  \fi
  \par\vspace{3pt}%
}